At the original threshold $e_i/d_i>1$, the seven-SNR reduction in
failure rate versus in-loop OptShrink is
4.13
percentage points, with paired interval
$[3.051,5.283]$.
At $-10$~dB alone the interval is
$[21.354,36.979]$
percentage points. Learned seed-specific rates there are
5.47\%, 5.47\% and 7.81\%, versus 35.16\% for OptShrink.

\begin{table}[!t]
\caption{Post-hoc error-threshold sensitivity on the stronger bank.
Rates average seven SNR cells; positive differences favor learning.}
\label{tab:tails}\centering\footnotesize\setlength{\tabcolsep}{3pt}
\begin{tabular}{@{}rrrr@{}}\toprule
NMSE & Learned & In-loop & OptShrink minus learned\\
threshold & (\%) & OptShrink (\%) & 95\% interval (points)\\\midrule
$0.5$ & 17.00 & 15.85 & $[-1.935,-0.409]$\\
$1$ & 0.89 & 5.02 & $[3.051,5.283]$\\
$2$ & 0.00 & 0.00 & ---\\
\bottomrule\end{tabular}
\end{table}
The conclusion depends on the error threshold (Table~\ref{tab:tails}).
At threshold 0.5, OptShrink has fewer exceedances; at threshold 2,
neither method has any observed exceedance. Zero observed events do not
establish zero population risk. These exploratory, pointwise intervals
retain paired worlds and fixed models and are not adjusted for multiple
thresholds. Learning therefore does not dominate the whole error
distribution. The original severe-error definition and primary NMSE
metric are unchanged.
